\section*{Parity certificates (author source pp.12,24)}
A parity certificate for an input $x\in\bits^n$ to a function $f$ is a set $A$ of linear forms such that $f$ is constant on the affine subspace defined by $A=A(x)$. A parity certificate complexity $C_\oplus(f,x)$ of a function $f$ at input $x$ is the size (number of linear forms) of the smallest parity certificate for $x$. The parity certificate complexity of $f$ is defined as $C_\oplus(f):=\max_{x\in\bits^n}C_\oplus(f,x)$.
\section*{5.2 Minimum weight lemma}
A parity certificate $C$ is a system of affine equations over input variables. It can be represented by a matrix $A$ and a vector $b$ that define a linear subspace, i.e., $Ax=b$.
\begin{originalresult}{Definition 5.2}
The minimum weight of a parity certificate $C$ is the minimum Hamming weight of nonzero vectors in the row space of $A$.
\end{originalresult}
For example, one can consider a certificate $\{x+y=1,x+y+z=0\}$. The minimum weight of this certificate is equal to 1, and this corresponds to the equation $z=1$. On the other hand, the minimum weight of a certificate $\{x+y=1,x+z=0\}$ is equal to 2.
